% Generated draft: requires mathematical and source review.
\documentclass[11pt]{article}
\usepackage[margin=25mm]{geometry}
\usepackage{fontspec}
\setmainfont{Latin Modern Roman}
\usepackage{amsmath,amssymb,mathtools,unicode-math}
\setmathfont{Latin Modern Math}
\usepackage{enumitem,needspace,xurl,hyperref}
\setlength{\parindent}{0pt}
\setlength{\parskip}{5pt}
\setlength{\emergencystretch}{3em}
\newcommand{\N}{\mathbb{N}}
\newcommand{\Z}{\mathbb{Z}}
\newcommand{\Q}{\mathbb{Q}}
\newcommand{\R}{\mathbb{R}}
\newcommand{\C}{\mathbb{C}}
\newcommand{\F}{\mathbb{F}}
\newcommand{\A}{\mathbb{A}}
\newcommand{\PP}{\mathbb{P}}
\newcommand{\E}{\mathbb{E}}
\newcommand{\eps}{\varepsilon}
\newcommand{\dd}{\,\mathrm{d}}
\begin{document}
\section*{UnsolvedMath ID: 30002089}
\textbf{Source problem code:} catalogue entry
\subsection{Definitions and mathematical statement}
A homological complex consists of objects \(C_i\) and differentials \(d_i:C_i\to C_{i-1}\), indexed by \(i\in\Z\), in a suitable additive setting, with \(d_{i-1}\circ d_i=0\). When the terms have classes in an appropriate additive group and the alternating sum is defined, its Euler characteristic is
\[
\chi(C_\bullet)=\sum_{i\in\Z}(-1)^i[C_i].
\]
The source does not specify the category of terms, coefficient system, additional gradings, or conditions making this sum meaningful.

Write \(B_p(D)\) provisionally for the parity bracket of a diagram \(D\) in the intended setting. This notation names the invariant in the supplied question; its definition, diagram domain, parity axioms, allowed moves, and value group are not included in the supplied record.

The requested research constructions are:
\begin{enumerate}
\item Construct a homological complex \(C_\bullet(D)\) categorifying the parity bracket, so that, for every diagram \(D\) in the intended domain,
\[
\chi\bigl(C_\bullet(D)\bigr)=B_p(D).
\]
The construction must specify the common target in which this equality is interpreted and the invariant homological information associated with the underlying object. The record does not specify the required notion of equivalence of complexes or the precise invariance requirement.
\item Generalize parity theory and an associated bracket to higher-dimensional objects, with surfaces in \(4\)-space explicitly proposed as an example. Specify the objects, their diagrams or presentations, the relevant equivalence moves, suitable parity data, and a bracket invariant compatible with that equivalence. The record does not specify the ambient category, permitted singularities, boundary conditions, or whether an extension must recover a particular existing bracket under a specified specialization.
\end{enumerate}
These are two construction questions. The displayed categorification identity is schematic until the missing domain, parity bracket, and Euler-characteristic target are fixed; the supplied record does not determine a unique fully specified mathematical formulation.
\subsection{Short English statement}
Construct a homological complex whose Euler characteristic recovers the parity bracket, and develop parity theory and a bracket for higher-dimensional objects such as surfaces in \(4\)-space. The supplied record leaves the underlying diagram theory, algebraic target, invariance conditions, and surface category unspecified, so this is a statement-only research-direction draft requiring mathematical review.
\subsection{Sources}
\begin{itemize}
\item Supplied local UnsolvedMath registry record, permanent ID \(30002089\), source problem code \texttt{catalogue entry}, titled \emph{Categorification and Higher-Dimensional Extensions of Parity Brackets}, in the Oberwolfach Reports Open Problems collection. No individual proposer is supplied.
\item Supplied report citation: \emph{Invariants in Low-Dimensional Topology and Knot Theory} (2013), Oberwolfach Reports, pp.~1687--1758, DOI \texttt{10.4171/owr/2012/28}. \url{https://doi.org/10.4171/owr/2012/28}. The registry separately records proposed year 2012. Both dating details are retained without resolving their relationship.
\item Supplied literature reference: Manturov, \emph{Free Knots and Parity}. \url{https://arxiv.org/abs/0912.5348}. This is a supplied background reference; the record does not attribute the present questions to an individual author.
\item Supplied literature reference labeled ``original Oberwolfach report'': \url{https://doi.org/10.4171/owr/2012/28}.
\item Literature assessment dated \texttt{2026-08-21}: the supplied curation labels the questions open and reports that no general construction resolving these research directions was identified. This is supplied dated curation derived from the OWR material, not an independently verified assessment; time-sensitive status and cited claims require independent verification.
\item The registry labels the statement ``exact'' and reports that its mathematical content was checked against the official report, with notation, context, and question form normalized without changing the source-level claim. That verification statement is supplied provenance and has not been independently checked for this draft.
\end{itemize}
\end{document}
